\documentclass[12pt]{article}
\usepackage{amsmath,amssymb,amsthm}
\usepackage{graphicx}
\usepackage{tikz}
\usetikzlibrary{calc}
\usepackage{booktabs}
\usepackage[shortlabels]{enumitem}
\setlist{itemsep=1pt,parsep=0pt,topsep=3pt}
\setlength{\emergencystretch}{2em}
\usepackage[colorlinks=true,linkcolor=blue!45!black,citecolor=blue!45!black,urlcolor=blue!45!black]{hyperref}
\hypersetup{pdftitle={Maximum Blackouts That Keep Every Lattice Rectangle Identifiable},pdfauthor={Arjun Pemmasani}}

\definecolor{gridcol}{gray}{0.86}
\definecolor{keptcol}{RGB}{226,236,247}
\definecolor{keptline}{gray}{0.35}
\definecolor{hidcol}{gray}{0.62}
\definecolor{pencol}{gray}{0.55}
\definecolor{cyccol}{RGB}{31,94,168}
\definecolor{rectA}{RGB}{0,114,178}
\definecolor{rectB}{RGB}{213,94,0}
\definecolor{rectC}{RGB}{0,158,115}
\definecolor{excesscol}{RGB}{253,214,160}
\tikzset{
  lab/.style={font=\scriptsize,text=black!55},
  capt/.style={font=\small},
  tiny/.style={font=\scriptsize},
  rowv/.style={draw,rectangle,fill=white,inner sep=1.5pt,minimum size=13pt,font=\scriptsize},
  colv/.style={draw,circle,fill=white,inner sep=0.5pt,minimum size=14pt,font=\scriptsize},
}

\theoremstyle{plain}
\newtheorem{theorem}{Theorem}
\newtheorem{lemma}[theorem]{Lemma}
\newtheorem{proposition}[theorem]{Proposition}
\newtheorem{corollary}[theorem]{Corollary}
\newtheorem{conjecture}[theorem]{Conjecture}
\newtheorem*{thmA}{Theorem A}
\newtheorem*{thmB}{Theorem B}
\newtheorem*{thmC}{Theorem C}
\newtheorem*{thmD}{Theorem D}
\newtheorem*{thmE}{Theorem E}
\theoremstyle{definition}
\newtheorem{definition}[theorem]{Definition}
\newtheorem{example}[theorem]{Example}
\theoremstyle{remark}
\newtheorem{remark}[theorem]{Remark}

\newcommand{\seqnum}[1]{\href{https://oeis.org/#1}{\rm \underline{#1}}}
\newcommand{\pres}{P_S}
\newcommand{\kA}{\mathsf{A}}
\newcommand{\kB}{\mathsf{B}}
\newcommand{\kC}{\mathsf{C}}
\newcommand{\kE}{\mathsf{E}}

\title{Maximum Blackouts That Keep Every Lattice Rectangle Identifiable}
\author{Arjun Pemmasani\\
Department of Mathematics\\
Harvey Mudd College\\
Claremont, CA 91711\\
USA\\
\texttt{apemmasani@hmc.edu}}
\date{}

\begin{document}
\maketitle

\begin{abstract}
Kagey's Problem 001 concerns a rectangular grid of points, some of them blacked out. A
rectangle with its four corners on grid points, tilted or not, shows only its white
corners. The problem asks how many points can be blacked out while every rectangle stays
determined by what it shows. We study the largest such blackout and the number of
blackouts of that size, the tables A397315 and A399580 of the On-Line Encyclopedia of
Integer Sequences. The main tool is a graph on the rows and columns whose short cycles
are axis-aligned confusions. With Mantel's theorem it bounds the maximum by the width
plus a quarter of the squared height. We find the maximum on up to five rows at every
width, and on six rows from width 15 on. On five rows from width 8 on, and on the
$6\times8$ grid, the maximum exceeds the sum of the two sides, against a guess we
recorded in the encyclopedia. On three rows we count the optima exactly, and on four
rows we find their growth rate. For every fixed number of rows the maximum minus the
width is eventually periodic. Some values rest on computer checks. Whether the bound is
attained on seven or more rows is open.
\end{abstract}

\input{sections/sec01_intro}
\input{sections/sec02_conv}
\input{sections/sec03_strip}
\input{sections/sec04_graph}
\input{sections/sec05_three}
\input{sections/sec06_four}
\input{sections/sec07_five}
\input{sections/sec08_height}
\input{sections/sec09_verify}
\input{sections/sec10_open}
\input{sections/sec11_ack}

\appendix
\input{sections/appB_threegf}

\begin{thebibliography}{99}
\small\setlength{\baselineskip}{11pt}\setlength{\itemsep}{0pt}\setlength{\parsep}{0pt}

\bibitem{alanko2011}
S. Alanko, S. Crevals, A. Isopoussu, P. R. J. \"Osterg\r{a}rd, and V. Pettersson,
Computing the domination number of grid graphs, \emph{Electron. J. Combin.}
\textbf{18} (2011), Paper P141.

\bibitem{alonspencer2016}
N. Alon and J. H. Spencer, \emph{The Probabilistic Method}, 4th ed., Wiley, 2016.

\bibitem{baccelli1992}
F. Baccelli, G. Cohen, G. J. Olsder, and J.-P. Quadrat, \emph{Synchronization and
Linearity: An Algebra for Discrete Event Systems}, Wiley, 1992.

\bibitem{bellitto2018}
T. Bellitto, Separating codes and traffic monitoring, \emph{Theoret. Comput. Sci.}
\textbf{717} (2018), 73--85.

\bibitem{butler2019}
S. Butler, J. Choi, K. Kim, and K. Seo, Enumerating multiplex juggling patterns,
\emph{J. Integer Sequences} \textbf{22} (2019), Article 19.1.7.

\bibitem{decaenszekely1997}
D. de Caen and L. A. Sz\'ekely, On dense bipartite graphs of girth eight and upper
bounds for certain configurations in planar point-line systems, \emph{J. Combin.
Theory Ser. A} \textbf{77} (1997), 268--278.

\bibitem{erdos1080}
P. Erd\H{o}s, Problem 1080, \emph{Erd\H{o}s Problems},
\url{https://www.erdosproblems.com/1080} (accessed October 2026).

\bibitem{euler2005}
R. Euler, The Fibonacci number of a grid graph and a new class of integer sequences,
\emph{J. Integer Sequences} \textbf{8} (2005), Article 05.2.6.

\bibitem{flajolet2009}
P. Flajolet and R. Sedgewick, \emph{Analytic Combinatorics}, Cambridge University
Press, 2009.

\bibitem{foucaud2015}
F. Foucaud, Decision and approximation complexity for identifying codes and
locating-dominating sets in restricted graph classes, \emph{J. Discrete Algorithms}
\textbf{31} (2015), 48--68.

\bibitem{glp2025}
J. Guilhot, E. Little, and J. Parkinson, Bounded weight functions on regular languages
and groups, arXiv preprint arXiv:2510.17061 [math.GR], 2025. Available at
\url{https://arxiv.org/abs/2510.17061}.

\bibitem{jiang2018}
M. Jiang, Periodicity of identifying codes in strips, \emph{Inform. Process. Lett.}
\textbf{135} (2018), 77--84.

\bibitem{kageyG4G14}
P. Kagey, Open problem collection, Problem 1, G4G14 gift exchange, 2020,
\url{https://www.gathering4gardner.org/g4g14gift/G4G14-PeterKagey-OpenishProblems.pdf}.

\bibitem{kagey001}
P. Kagey, Problem 001, Open Problem Collection,
\url{https://peterkagey.com/problems/001/} (accessed October 2026).

\bibitem{katona1966}
G. Katona, On separating systems of a finite set, \emph{J. Combin. Theory}
\textbf{1} (1966), 174--194.

\bibitem{mantel1907}
W. Mantel, Problem 28, \emph{Wiskundige Opgaven} \textbf{10} (1907), 60--61.

\bibitem{mertens2024}
S. Mertens, Domination polynomials of the grid, the cylinder, the torus, and the king
graph, \emph{J. Integer Sequences} \textbf{27} (2024), Article 24.7.6.

\bibitem{pulp}
S. Mitchell, M. O'Sullivan, and I. Dunning, PuLP: a linear programming toolkit for
Python, technical report, 2011, \url{https://optimization-online.org/2011/09/3178/}.
We use version 3.3.2.

\bibitem{neuwirth2001}
S. Neuwirth, The size of bipartite graphs with girth eight, arXiv preprint
arXiv:math/0102210 [math.CO], 2001. Available at
\url{https://arxiv.org/abs/math/0102210}.

\bibitem{oeis}
OEIS Foundation Inc., \emph{The On-Line Encyclopedia of Integer Sequences}, 2026,
\url{https://oeis.org}.

\bibitem{trureturing2026}
The Omega Institute, Prove the exact Three-Row blackout maximum, pull request 8084,
merged on 15 September 2026 into the branch \texttt{dev} as commit \texttt{ace272f},
\url{https://github.com/the-omega-institute/trureturing/pull/8084}.


\bibitem{repo}
A. Pemmasani, Kagey-Problems repository, folder \texttt{kagey-problems/problem001},
2026, \url{https://github.com/apemm/Kagey-Problems}.

\end{thebibliography}

\noindent 2020 {\it Mathematics Subject Classification}: Primary 05A15; Secondary
05C35, 05B40, 52C05.

\noindent \emph{Keywords:} blackout, lattice rectangle, test cover, girth, transfer matrix.

\noindent (Concerned with sequences
\seqnum{A001787},
\seqnum{A085582},
\seqnum{A289832},
\seqnum{A322284},
\seqnum{A397315}, and
\seqnum{A399580}.)

\end{document}
